%% !TEX root = manuscript.tex

\chapter{Organisation, compétences et prise de recul}
\minitoc
\newpage

Ce dernier chapitre quitte le récit des missions pour revenir sur la façon dont j'ai organisé
mon travail, sur les compétences que j'ai mobilisées et acquises, sur les difficultés
rencontrées, sur l'adéquation de ces missions avec ma formation et sur le positionnement RSE de
l'entreprise. Il se termine par mes perspectives et mon projet professionnel.

\section{Gestion de projet et organisation}

Mon organisation a été largement conditionnée par un fait que j'ai déjà signalé : le groupe ne
dispose pas d'une équipe data constituée. Je n'ai donc pas travaillé dans le cadre d'une méthode
Agile formalisée, avec ses sprints et ses rituels, comme c'est le cas dans une direction data
structurée. Mon organisation a été plus artisanale, et il me semble utile de le dire, parce que
cela a eu des conséquences concrètes.

Le calendrier de l'année a d'abord été imposé de l'extérieur, par le projet de changement d'ERP,
qui commandait les échéances de la reprise des données. Autour de ce fil principal, j'ai
fonctionné par jalons : un lot de test, une validation, un lot plus large. Cette progression par
paliers, née d'une contrainte technique sur la migration, m'a finalement servi de méthode de
gestion de projet sur toutes mes missions, y compris sur le Copilote Financier, où le cadrage
puis la restitution ont été validés avec la direction financière avant d'aller plus loin.

Sur le plan personnel, la difficulté a été de tenir ensemble deux rythmes très différents. D'un
côté, les projets de fond, qui demandent de longues plages de concentration. De l'autre, les
sollicitations du quotidien, une extraction qui ne remonte plus, un accès à ouvrir, un chiffre
à expliquer, qui arrivent sans prévenir et qui attendent une réponse rapide. Hiérarchiser mes
tâches et m'y tenir a demandé un effort constant, et je constate qu'il est facile de se laisser
emporter par l'urgence au détriment du fond. C'est une difficulté sur laquelle je reviens plus
loin.

Le second semestre a marqué une évolution nette de mes pratiques sur ce point. J'ai versionné
l'intégralité du Copilote Financier sous Git, tenu un journal de bord documentant chaque étape
et ses enseignements, et imposé au projet une contrainte de reproductibilité complète, tout
devant se reconstruire depuis les exports bruts par une seule commande. Ces trois pratiques
n'existaient pas dans mon travail du premier semestre, et je les considère comme un acquis
durable.

\section{Compétences mobilisées et acquises}

Sur le plan technique, j'ai consolidé et surtout élargi mes bases. J'ai beaucoup pratiqué les
interfaces de programmation, du simple appel de test au requêtage dynamique et à la modification
partielle au format JSON Patch, ainsi que le tableur employé non comme un outil de saisie mais
comme un moteur de transformation piloté par formules. J'ai également travaillé SQL, de la
requête d'extraction à la conception d'une base analytique avec ses dimensions, ses tables de
faits et ses vues métier. Sur le volet analyse, j'ai mis en pratique un ensemble de méthodes
vues en cours, de la régression linéaire aux modèles à effets aléatoires, en passant par les
séries temporelles, la classification, l'optimisation sous contrainte et le traitement du texte.
J'ai enfin découvert Streamlit et Git, que je ne pratiquais pas avant cette année.

Sur le plan méthodologique, l'acquis principal est la validation. J'ai appris à ne jamais livrer
un résultat sans avoir défini d'abord comment je saurais qu'il est juste, et surtout à
comprendre que cette question ne se pose pas de la même façon selon le livrable.

Sur le plan humain et professionnel, j'ai gagné en autonomie, ce qui n'était pas un choix mais
une nécessité compte tenu de la taille de l'équipe. J'ai surtout gagné en capacité à dialoguer
avec des profils non techniques, c'est-à-dire à traduire un besoin métier flou en règle codable,
et parfois à en discuter la faisabilité. J'ai également développé une pratique de la
rétro-ingénierie, en reprenant des systèmes dont les règles n'étaient documentées nulle part.

\section{L'intelligence artificielle dans ma pratique}

Il me semble important de consacrer une place à l'intelligence artificielle dans ce mémoire,
tant elle a marqué ma façon de travailler cette année.

Concrètement, elle m'a servi d'appui au quotidien, notamment pour monter rapidement en
compétence sur des outils que je découvrais, pour m'aider à écrire ou à corriger du code, pour
clarifier une notion vue en cours, ou encore pour structurer un document. Elle a agi comme un
accélérateur d'apprentissage, particulièrement précieux dans un contexte où je devais me former
seule à un environnement technique nouveau, sans équipe data autour de moi pour prendre le
relais.

Cet usage ne va toutefois pas sans réflexion. Une intelligence artificielle produit vite, mais
pas toujours juste, et ses réponses peuvent être plausibles tout en étant erronées. J'ai donc
fait attention à ne jamais prendre une sortie d'IA pour argent comptant, et à la vérifier
systématiquement, exactement comme je le ferais pour une donnée. L'outil me fait gagner du
temps, mais la responsabilité du résultat me revient entièrement.

Par souci de transparence, je précise enfin que l'IA m'a assistée dans la rédaction et la mise
en forme de ce mémoire, comme dans une partie de mon travail technique, toujours sous mon
contrôle et ma relecture.

\section{Retour d'expérience : difficultés et enseignements}

Cette année a été marquée par des difficultés qui se sont toutes révélées formatrices. Trois
enseignements me paraissent mériter d'être développés, parce qu'ils ne ressortent que de la
comparaison entre mes missions.

\subsection{Traduire avant de modéliser}

L'image que j'avais de ce métier en entrant en Master~2 était celle que donne le cursus : on
reçoit un jeu de données et une question, on choisit une méthode, on évalue, on livre un modèle.
Cette image n'est pas fausse, elle est simplement très minoritaire dans le temps réel. Sur les
quatre missions de l'année, une seule y correspond, et encore partiellement. Le poste que j'ai
réellement occupé se rapproche davantage de celui de gestionnaire d'applications~[1].

Ce qui a occupé le reste, c'est un travail de traduction dans les deux sens : traduire une gêne
exprimée par un service en question calculable, puis traduire un résultat statistique en
information sur laquelle quelqu'un peut décider. Cette traduction est difficile pour une raison
précise, c'est que les deux langues n'ont pas les mêmes catégories. Le service juridique parle
de commandes qui se suivent, alors que la base ne connaît que des lignes sans position dans une
séquence. La direction financière parle d'opérations qui dérivent, alors que les exports ne
contiennent qu'un budget et un engagé à une date. Chaque fois, il a fallu construire la
catégorie manquante, et ce sont ces constructions qui engagent le plus, parce qu'elles décident
de ce que l'entreprise pourra observer ensuite.

\subsection{Vérifier son travail ne veut pas dire la même chose selon le livrable}

C'est l'enseignement transversal qui m'a le plus servi, et je ne l'ai vu qu'en comparant mes
missions entre elles. Une reprise de données, un tableau de bord et un modèle statistique ne se
vérifient pas de la même manière, et la figure~\ref{fig:verification} récapitule les trois
manières que j'ai employées.

\begin{figure}[!htbp]
\centering
\includegraphics[width=0.85\textwidth]{fig_5_1_verification.png}
\caption{Les trois manières de vérifier un résultat, selon la nature du livrable}
\label{fig:verification}
\end{figure}

Sur une reprise de données, la question est de savoir si ce qui est arrivé correspond à ce qui
est parti, et j'y ai répondu par des paliers de validation, des compteurs de contrôle et le
recoupement de deux sources. Sur un modèle, la question est de savoir si le résultat tient hors
de l'échantillon d'apprentissage, et j'y ai répondu par un échantillon test, une comparaison à
une référence naïve et un contrôle des fuites. Sur un livrable destiné à un utilisateur, la
question est enfin de savoir s'il s'en sert, et pour la bonne chose, et cela ne se vérifie que
par la validation du service et son retour à l'usage.

J'ai appris à mes dépens qu'appliquer la mauvaise vérification ne protège de rien. Les
compteurs de la reprise des tiers tombaient juste, mais ils n'auraient rien dit d'une clé
d'unicité mal choisie. Ce sont les échanges avec les équipes commerciales, et non les
compteurs, qui ont validé ces clés.

\subsection{Documenter fait partie du travail technique}

J'ai longtemps considéré la documentation comme une formalité qui suit le travail. Deux épisodes
m'ont fait changer d'avis. Le premier est la note interne sur la séquence d'injection dans
l'ERP, rédigée au premier semestre, qui m'a servi deux mois plus tard sur la reprise des tiers
et qui aurait servi à quelqu'un d'autre si je n'avais pas été là. Le second est la note de
procédure de l'import des tiers, qui énonce les manipulations à ne jamais faire avant même
d'expliquer comment utiliser le fichier.

Dans un contexte où je suis la seule à connaître ces règles, ne pas les écrire revient à créer
une dépendance à ma présence, ce qui est un défaut du livrable et non une preuve de sa valeur.

\subsection{Bilan des livrables}

Le tableau suivant récapitule ce qui a été livré et se trouve aujourd'hui en service.

\begin{longtable}[]{@{}
  >{\raggedright\arraybackslash}p{(\columnwidth - 2\tabcolsep) * \real{0.6329}}
  >{\raggedright\arraybackslash}p{(\columnwidth - 2\tabcolsep) * \real{0.3671}}@{}}
\toprule\noalign{}
\begin{minipage}[b]{\linewidth}\raggedright
Livrable
\end{minipage} & \begin{minipage}[b]{\linewidth}\raggedright
Volume ou périmètre
\end{minipage} \\
\midrule\noalign{}
\endhead
\bottomrule\noalign{}
\endlastfoot
Opérations créées dans SPO, hors budgets détaillés & 283 \\
Tiers créés dans SPO, avec leurs objets rattachés & 1~434 \\
Adresses, courriels et téléphones importés & 553, 469 et 538 \\
Représentants des couples et des sociétés & 1~871 \\
Coordonnées ajoutées par requête API & 872 requêtes, toutes passées \\
Documentation interne produite & Séquence d'injection, procédure d'import des tiers \\
Reportings maintenus et faits évoluer & Sur l'ensemble des services demandeurs \\
Postes équipés du SSO & Agence de Perpignan et service financier \\
Base de données analytique & Trois dimensions, quatre tables de faits, trois vues métier \\
Application d'aide à la décision livrée & Cinq pages, 147 opérations de promotion suivies \\
\end{longtable}

\section{Adéquation avec le Master MIASHS}

Le Master MIASHS, Mathématiques et Informatique Appliquées aux Sciences Humaines et Sociales,
consiste à mobiliser les mathématiques, les statistiques et l'informatique pour analyser des
phénomènes humains et sociaux. Mon alternance rejoint cette formation sur plusieurs points, mais
elle appelle aussi une nuance que je préfère poser d'emblée.

Sur le plan méthodologique, le Copilote Financier a mobilisé très directement les enseignements
du Master. La régression linéaire et la régularisation, la classification et ses métriques,
l'analyse exploratoire et ses contrôles, les séries temporelles, l'optimisation sous contrainte
et le traitement du texte sont autant de notions vues en cours que j'ai appliquées sur des
données réelles, avec les difficultés que cela suppose. Le volet informatique de la formation
m'a été tout aussi utile, pour l'écriture des requêtes SQL comme pour la structuration d'un
projet Python reproductible.

Au-delà des méthodes, l'objet d'étude lui-même rejoint le cursus. Derrière chaque ligne de mes
tables, il y a des ménages qui achètent un logement, des projets qui échouent faute de
financement, et un marché qui réagit au coût du crédit. L'axe B de mon projet est, au fond, une
étude de la sensibilité d'une demande à une variable macroéconomique, ce qui relève pleinement
des sciences sociales quantitatives.

La nuance est la suivante. Sur les quatre missions de l'année, trois relèvent de l'ingénierie
des données et de la gestion d'applications, et non de la modélisation. Le Master accorde une
place importante à la data science au sens strict, et mon alternance m'a surtout placée en
amont de ces usages. Je ne le regrette pas, parce que c'est précisément ce décalage qui m'a fait
comprendre ce dont dépend une modélisation ; mais je dois reconnaître que la part de mon année
consacrée à la modélisation est restée minoritaire.

\section{Positionnement RSE et propositions}

Ma position d'alternante dans une fonction support ne m'a pas donné de visibilité sur la
politique RSE du groupe dans son ensemble, et je me garderai donc d'en dresser un bilan. Je peux
en revanche décrire ce que j'ai observé et ce à quoi mon travail a contribué, à mon échelle.

Sur le plan social, l'entreprise m'a offert un cadre d'apprentissage réel : j'ai été confiée à
des chantiers structurants, avec l'autonomie et le droit à l'erreur que cela suppose, ce qui est
loin d'être acquis pour une alternante. Le fait d'avoir été la seule personne sur ces sujets a
été exigeant, mais aussi profondément formateur.

Sur le plan environnemental, ma contribution est indirecte et je ne veux pas la surestimer. La
suppression de la double saisie entre systèmes, la réduction du nombre de fiches clients
dupliquées et la substitution d'une chaîne de traitement automatisée à des centaines de saisies
manuelles réduisent surtout du temps de travail et des erreurs, pas une empreinte carbone. Le
seul effet mesurable dans ce registre concerne le Copilote Financier, dont l'architecture
calcule une seule fois, dans une vue, ce que des tableurs recalculaient à chaque ouverture.

Je formulerais deux propositions, qui sont de même nature que mes recommandations techniques.
La première est l'historisation des exports mensuels dans le nouvel ERP, qui donnerait à
l'entreprise une mémoire de sa propre activité, aujourd'hui perdue à chaque nouvel export. La
seconde est la fiabilisation de la saisie des motifs de désistement, dont l'analyse a montré
qu'elle rend, dans certaines agences, une part importante de l'information inexploitable. Ces
deux chantiers ne coûtent presque rien et conditionnent tout ce qui pourra être analysé dans les
années à venir.

\section{Perspectives et projet professionnel}

Ces deux années ont déplacé mon projet. J'y suis entrée en visant la data science au sens
strict, c'est-à-dire la modélisation prédictive. J'en sors en visant un positionnement plus
large : construire et fiabiliser les chaînes de données qui rendent la décision possible, puis
les outiller.

Ce positionnement croise trois domaines que l'année m'a fait pratiquer ensemble.
L'\textbf{ingénierie des données} d'abord, sur laquelle j'ai le plus progressé, avec les
reprises, les migrations, les interfaces de programmation, la déduplication et la modélisation
d'une base analytique. L'\textbf{analyse et la modélisation} ensuite, entendues comme le choix
d'une méthode proportionnée à la donnée disponible et sa validation honnête. La
\textbf{restitution et le dialogue avec les métiers} enfin, c'est-à-dire la traduction d'un
besoin en spécification, puis d'un résultat en écran. Ces trois domaines se tiennent : un profil
qui ne tiendrait que le premier produirait des chaînes de données que personne n'interroge ; un
profil qui ne tiendrait que le deuxième modéliserait des données dont il ignore la provenance.

À court terme, je vise un poste d'ingénieure données orientée aide à la décision. Deux contextes
m'intéressent, pour des raisons différentes. Une entreprise de taille intermédiaire sans équipe
data constituée, où le poste couvre toute la chaîne, comme cela a été le cas cette année. Ou une
direction data structurée, où je travaillerais enfin à plusieurs sur un même chantier, ce qui me
manque et ce dont j'ai besoin pour progresser sur les pratiques collectives : la revue de code,
les conventions partagées et la répartition du travail.

À plus court terme encore, plusieurs chantiers restent ouverts chez Angelotti, et je souhaite
les mener à terme : achever la reprise des historiques budgétaires, construire la
synchronisation permanente entre l'ERP et le CRM, et accompagner l'adoption du Copilote
Financier par la direction financière, dont l'usage réel sera le meilleur juge.
